\section {Текст задания}

По варианту, выданному преподавателем, составить и выполнить запросы к базе данных "Учебный процесс".

Команда для подключения к базе данных ucheb:

psql -h pg -d ucheb

Отчёт по лабораторной работе должен содержать:
\begin{itemize}
	\item Текст задания.
	\item Реализацию запросов на SQL.
	\item Выводы по работе.
\end{itemize}

Темы для подготовки к защите лабораторной работы:
\begin{itemize}
	\item SQL
	\item Соединение таблиц
	\item Подзапросы
	\item Представления
	\item Последовательности
\end{itemize}

Составить запросы на языке SQL (пункты 1-7).
\begin{enumerate}
	\item Сделать запрос для получения атрибутов из указанных таблиц, применив фильтры по указанным условиям:
	
	Таблицы: Н\_ЛЮДИ, Н\_СЕССИЯ.
	
	Вывести атрибуты: Н\_ЛЮДИ.ИМЯ, Н\_СЕССИЯ.УЧГОД.
	
	Фильтры (AND):
	\begin{enumerate}
		\item Н\_ЛЮДИ.ФАМИЛИЯ > Ёлкин.
		\item Н\_СЕССИЯ.ИД < 14.
		\item Н\_СЕССИЯ.ИД = 32199.
	\end{enumerate}
	Вид соединения: RIGHT JOIN.
	\item Сделать запрос для получения атрибутов из указанных таблиц, применив фильтры по указанным условиям:
	
	Таблицы: Н\_ЛЮДИ, Н\_ВЕДОМОСТИ, Н\_СЕССИЯ.
	
	Вывести атрибуты: Н\_ЛЮДИ.ИМЯ, Н\_ВЕДОМОСТИ.ЧЛВК\_ИД, Н\_СЕССИЯ.УЧГОД.
	
	Фильтры (AND):
	\begin{enumerate}
		\item Н\_ЛЮДИ.ИД > 152862.
		\item Н\_ВЕДОМОСТИ.ДАТА < 2010-06-18.
	\end{enumerate}
	Вид соединения: LEFT JOIN.
	\item Вывести число студентов вечерней формы обучения, которые младше 20 лет.
	Ответ должен содержать только одно число.
	\item В таблице Н\_ГРУППЫ\_ПЛАНОВ найти номера планов, по которым обучается (обучалось) менее 2 групп на кафедре вычислительной техники.
	Для реализации использовать соединение таблиц.
	\item Выведите таблицу со средними оценками студентов группы 4100 (Номер, ФИО, Ср\_оценка), у которых средняя оценка равна минимальной оценк(е|и) в группе 1100.
	\item Получить список студентов, зачисленных после первого сентября 2012 года на первый курс заочной формы обучения (специальность: Программная инженерия). В результат включить:
	номер группы;
	номер, фамилию, имя и отчество студента;
	номер и состояние пункта приказа;
	Для реализации использовать подзапрос с EXISTS.
	\item Сформировать запрос для получения числа в СПбГУ ИТМО отличников.
\end{enumerate}






